\section{Introduction}\label{sec:intro}

Strikes, natural disasters, armed conflicts, and the closure of a port, bridge, or strait can cut critical routes of a supply chain and sharply reduce the amount of goods a logistics network can deliver~\cite{supplychain_resilience}. Identifying the routes whose loss costs the most capacity tells planners where to invest in protection and redundancy. Operations research models this question as the \emph{maximum flow network interdiction problem} (MFNIP)~\cite{wood1993,smith2020}. The network has a source, a sink, and a capacity on each arc. An interdictor with a limited budget removes arcs, and a network user, the follower, then sends as much flow as possible from the source to the sink through the remaining arcs. The interdictor seeks the removal plan that minimizes this maximum flow. The problem is bilevel because evaluating even a single plan requires solving an optimization problem, namely the follower's maximum flow. A maximum flow can be computed in polynomial time, but the general problem of choosing the best plan is NP-hard~\cite{wood1993}. Figure~\ref{fig:example} in Section~\ref{sec:method} shows a small instance.

This paper asks what quantum computers can offer for a problem of this kind, and three facts constrain the answer. First, exact classical methods perform well in practice. Wood~\cite{wood1993} reformulated the MFNIP as a single-level mixed-integer linear program (MILP) by replacing the follower's maximum flow with the equivalent minimum cut (Section~\ref{sec:bg-mfnip}). In our scaling study, a modern MILP solver certified optimal plans for networks with thousands of nodes within 20 seconds (Section~\ref{sec:results-classical}), whereas the networks that fit on current quantum devices have a few dozen arcs and are trivial for classical solvers. Second, current quantum processors belong to the noisy intermediate-scale quantum (NISQ) era~\cite{preskill2018}. They have few qubits, limited connectivity, and noise, so they can only run shallow circuits. Third, the quantum optimizers most often applied to combinatorial problems, quantum annealing~\cite{kadowaki1998} and the quantum approximate optimization algorithm (QAOA)~\cite{farhi2014}, require a quadratic unconstrained binary optimization (QUBO) model~\cite{lucas2014,glover2019}. For the MFNIP, such a model must encode the node partition, the cut, and the budget as penalty terms, with slack variables for the inequalities. This encoding adds qubits, lets the device return infeasible plans, and comes with no performance guarantee. A quantum method for the MFNIP should therefore not be presented as a competitor of the MILP at the current scale. What it can offer is a provable speedup over a well-defined classical reference, together with an honest account of its cost in quantum resources.

The question that guides this work is \emph{what speedup a quantum search algorithm can provably achieve over classical exhaustive search for the MFNIP, at what cost in qubits and gates, and how it compares with the classical method used in practice.} Our objective is to design, analyze, and verify a quantum interface that applies the quantum minimum-finding algorithm of Dürr and Høyer~\cite{durr1996}, a Grover-type method~\cite{grover1996,bbht1998}, to the MFNIP without losing its quadratic speedup. The quadratic speedup of Grover search over brute force is well known, but applying it to a bilevel problem raises two difficulties. The oracle must compute the follower's maximum flow for every plan in superposition, and the budget must be enforced without enlarging the search space.

The work proceeds in six steps. We (i) restrict the problem to its cardinality variant, in which every arc costs one unit and at most $k$ arcs are removed, and compare the number $N$ of feasible plans with the size of the hypercube; (ii) design the quantum interface, which comprises the encoding, a state preparation that enforces the budget, a threshold predicate on the maximum flow, and the corresponding phase oracle; (iii) prove the query complexity of minimum finding with this interface and show that it is optimal; (iv) generate instances of increasing size with fixed seeds and validate each simulator against an exact reference; (v) measure in simulation how the cost of finding the optimum grows with $N$ for the quantum algorithm and for classical rivals, using the MILP both as the reference optimum and as the practical baseline; and (vi) count the qubits and gates of the interface as functions of the network size and run amplitude amplification on an IBM quantum processor.

We state two refutable hypotheses, fixed before the experiments (Section~\ref{sec:method-hypotheses}). Success always means finding an exactly optimal plan certified by the MILP. Hypothesis H1 concerns simulation. On 30 layered networks with 6 to 27 nodes, 8 to 70 arcs, and $N$ between 9 and 974{,}121, the mean cost of minimum finding until an optimal plan is first measured grows as $N^{\beta}$ with $\beta\in[0.4,0.6]$, whereas a classical random-order search grows with an exponent close to one. Because our simulator samples from the exact amplification probabilities, H1 checks the correctness of the implementation rather than the behavior of a physical oracle. Hypothesis H2 concerns hardware. On an IBM processor and for a network with three nodes and three arcs ($k=1$, $N=4$), the frequency of the optimal plan after the optimal number of Grover iterations is significantly greater than the uniform baseline $1/N$ (one-sided binomial test at level $0.05$). Since this baseline is weak, we also compare the hardware frequencies with the theoretical probabilities and with a noisy simulator in an exploratory analysis.

\paragraph{Contributions} This paper makes four contributions.
\begin{enumerate}
\item[C1.] A \emph{quantum interface} for the MFNIP (Section~\ref{sec:interface}). It uses one qubit per arc and a state preparation based on Dicke states~\cite{bartschi2019} that enforces the cardinality budget, so that the search covers exactly the $N=\sum_{i\le k}\binom{m}{i}$ feasible plans. A threshold predicate on the maximum flow, its phase oracle, and a classical threshold loop complete the interface.
\item[C2.] A \emph{provable speedup} (Proposition~\ref{prop:advantage} in Section~\ref{sec:advantage}). With $\bigO(\sqrt{N})$ applications of polynomial-size circuits, the algorithm finds an optimal plan with probability at least $1/2$, whereas any classical algorithm that only evaluates candidate plans needs $\Omega(N)$ flow evaluations. The bound is optimal in the query model.
\item[C3.] A \emph{verification} of the interface (Section~\ref{sec:results}). It comprises the growth of the cost to the optimum in a query-level simulation, the qubit and gate counts of the state preparation and of an arithmetic oracle, and the amplification of the optimal plan on ideal and noisy simulators and on an IBM quantum processor.
\item[C4.] A \emph{delimitation} of the speedup (Sections~\ref{sec:advantage} and~\ref{sec:discussion}). The speedup is quadratic, asymptotic, and relative to exhaustive search over the feasible plans, and on every instance we studied the MILP is orders of magnitude faster than both searches.
\end{enumerate}
A systematic literature search (Section~\ref{sec:rel-novelty}) found quantum approaches to other interdiction problems~\cite{peddireddy2026,sutcliffe2026}, but none for the MFNIP and none in which the state preparation enforces the interdiction budget.

\paragraph{Organization} Section~\ref{sec:background} reviews the background and Section~\ref{sec:related} the related work. Section~\ref{sec:interface} presents the interface, and Section~\ref{sec:advantage} proves the speedup. Section~\ref{sec:method} describes the experiments, and Section~\ref{sec:results} reports the results. Section~\ref{sec:discussion} discusses the results and their limitations, and Section~\ref{sec:conclusion} concludes.
